\begin{helper}
Let \(G\) be an \((n,d,\lambda)\)-loop graph. If \(f,g\colon V(G)\to\mathbb R\) satisfy
\[
  \sum_{v\in V(G)}f(v)=0,
  \qquad
  \sum_{v\in V(G)}g(v)=0,
\]
then
\[
  \left|\sum_{v\in V(G)} f(v)(A_Gg)(v)\right|
  \le
  \lambda
  \left(\sqrt{\sum_{v\in V(G)} f(v)^2}\right)
  \left(\sqrt{\sum_{v\in V(G)} g(v)^2}\right).
\]
\end{helper}
\begin{proof}
Use the standard real inner product on functions \(V(G)\to\mathbb R\),
\[
  \langle r,s\rangle=\sum_{v\in V(G)}r(v)s(v),
\]
and write \(T=A_G\) for the adjacency action from \(\noderef{LoopGraphAdjacencyAction}\).

First \(T\) is self-adjoint. Indeed, by \(\noderef{LoopGraphAdjacencyAction}\),
\[
\begin{aligned}
  \langle r,Ts\rangle
  &=\sum_v r(v)\sum_w {\bf 1}_{vw\in E(G)}s(w)  \\
  &=\sum_{v,w}{\bf 1}_{vw\in E(G)}r(v)s(w).
\end{aligned}
\]
The symmetry clause in \(\noderef{LoopGraphNdLambda}\), interpreted through \(\noderef{LoopGraphSymmetric}\), says \(vw\in E(G)\) if and only if \(wv\in E(G)\). Swapping \(v\) and \(w\) in the finite double sum therefore gives
\[
  \langle r,Ts\rangle=\langle Tr,s\rangle.
\]

Let
\[
  W=\left\{h:V(G)\to\mathbb R:\sum_v h(v)=0\right\}.
\]
The hypotheses say \(f,g\in W\). By the degree clause in \(\noderef{LoopGraphNdLambda}\), interpreted through \(\noderef{LoopGraphDegree}\), the adjacency action sends the constant-one function to \(d\) times the constant-one function. Thus for any \(h\in W\),
\[
  \sum_v (Th)(v)=\langle {\bf 1},Th\rangle
  =\langle T{\bf 1},h\rangle
  =d\langle {\bf 1},h\rangle=0.
\]
So \(W\) is \(T\)-invariant. The restriction of \(T\) to the finite-dimensional real inner product space \(W\) is still self-adjoint. By the finite-dimensional spectral theorem for self-adjoint operators, \(W\) has an orthonormal basis \(e_1,\dots,e_m\) consisting of eigenvectors of this restriction, say \(Te_i=\mu_i e_i\).

Each \(e_i\) has norm \(1\), so it is nonzero; it also lies in \(W\), so its coordinate sum is \(0\). Therefore every \(\mu_i\) is a non-principal adjacency eigenvalue in the sense of \(\noderef{LoopGraphNonprincipalEigenvalue}\). The eigenvalue-bound clause of \(\noderef{LoopGraphNdLambda}\) gives
\[
  |\mu_i|\le \lambda
  \qquad\text{for every }i.
\]
The final clause of \(\noderef{LoopGraphNdLambda}\) gives \(\lambda\ge 0\).

Expand \(f=\sum_i a_i e_i\) and \(g=\sum_i b_i e_i\). Since the \(e_i\) are orthonormal and \(Te_i=\mu_i e_i\),
\[
  \langle f,Tg\rangle
  =
  \sum_i \mu_i a_i b_i.
\]
Hence
\[
  |\langle f,Tg\rangle|
  \le
  \sum_i |\mu_i|\,|a_i|\,|b_i|
  \le
  \lambda\sum_i |a_i|\,|b_i|.
\]
By Cauchy's inequality,
\[
  \sum_i |a_i|\,|b_i|
  \le
  \sqrt{\sum_i a_i^2}\sqrt{\sum_i b_i^2}.
\]
Finally, Parseval's identity for the orthonormal basis gives
\[
  \sum_i a_i^2=\langle f,f\rangle=\sum_v f(v)^2,
  \qquad
  \sum_i b_i^2=\langle g,g\rangle=\sum_v g(v)^2.
\]
Substituting these identities into the preceding bound proves the claimed inequality.
\end{proof}
